With the increasing number of undergraduate theses, course papers, and various academic writing tasks in universities, the standardization of paper typesetting and the efficiency of editing have become important concerns for students and researchers. Although traditional word processing tools are intuitive to use, they still require considerable manual adjustment when dealing with complex formulas, figure and table numbering, cross-references, reference management, and thesis template adaptation. LaTeX has strong advantages in academic typesetting, but its command syntax, environment structures, compilation configuration, and error log analysis are not friendly to beginners. To address problems such as the high cost of memorizing LaTeX code, complex template usage, and difficulty in locating compilation errors, this thesis designs and implements an intelligent academic paper typesetting system based on JavaWeb and LaTeX, providing online, component-based, and intelligent support for academic writing and typesetting.

The system adopts a front-end and back-end separated architecture. The front end is developed based on Vue 3 and Vite to implement page display and user interaction, while the back end uses Spring Boot, MyBatis, and MySQL to complete business logic processing and data persistence. Centered on the academic paper writing process, the system implements functions such as user registration and login, paper project management, online LaTeX editing, template management, paper compilation, PDF preview, file export, AI dialogue assistance, and compilation error analysis. Users can create paper projects directly in the browser, select or import paper templates, edit paper content online, and generate PDF documents through different compilation engines.

The core design of this system lies in the introduction of a component library-based drag-and-drop insertion mechanism oriented toward paper structures. The system encapsulates common LaTeX paper structures, such as chapters, images, tables, formulas, lists, and symbols, into visual components. Users do not need to memorize complete LaTeX commands and environment syntax. Instead, they only need to select the required component from the component library and drag it into the editing area, after which the system automatically generates the corresponding LaTeX code framework. This approach lowers the threshold for using LaTeX syntax, reduces common problems such as command spelling errors, missing brackets, and unclosed environments, and improves paper editing efficiency.

During implementation, the system also optimizes problems such as numerous LaTeX template dependency files and difficult-to-understand compilation error messages. The system supports zip template package import and preserves the original directory structure as much as possible during template parsing and compilation, reducing compilation failures caused by missing dependency files such as .cls, .sty, and .bib. Meanwhile, the system introduces AI-assisted functions, which can provide cause analysis and modification suggestions based on compilation error logs, helping users quickly locate problems. Test results show that the system can complete major tasks such as paper project management, component drag-and-drop insertion, template loading, online editing, paper compilation, PDF preview, file export, and AI error analysis. The system runs relatively stably, basically meets the needs of online academic paper typesetting, and has certain practical application value.